\section{Ajuste del sistema: métricas, hardware y estrategias de servicio}

\subsection{Los distintos objetivos de optimización de TTFT y TPS}
Desde la perspectiva del servicio, Prefill y Decode corresponden a dos tipos de métricas de experiencia de usuario completamente distintos. Prefill afecta más directamente la latencia del primer token (TTFT), que determina «si el usuario siente que el sistema responde con suficiente rapidez»; Decode afecta más directamente la velocidad de salida por segundo (TPS), que determina la fluidez percibida en respuestas largas. Solo al entender esta diferencia se pueden configurar correctamente el procesamiento por lotes, el grado de concurrencia y el presupuesto de hardware.

\begin{table}[H]
\centering
\caption{Los dos tipos de objetivos más comunes en el servicio de inferencia}
\begin{tabular}{p{3.2cm}p{4.8cm}p{4.5cm}}
\toprule
Métrica & Etapa que más la afecta & Medios de optimización típicos \\
\midrule
TTFT & Prefill & Aumentar el rendimiento de GEMM, reducir la redundancia del prompt, compilar el kernel de antemano \\
TPS & Decode & Mejorar el ancho de banda de memoria, optimizar la KV cache, aumentar el procesamiento por lotes de decode \\
Relación costo/rendimiento & Efecto conjunto de ambas & Ajustar la concurrencia, el tamaño de lote, la longitud del contexto y la estrategia de cuantización \\
\bottomrule
\end{tabular}
\end{table}

\subsection{Por qué el mismo hardware muestra «dos caras»}
Una GPU puede comportarse como un «monstruo de cómputo» durante la etapa de Prefill, y sin embargo, en la etapa de Decode, parecer un «transportista limitado por el ancho de banda». Esto no significa que el hardware se haya vuelto peor de repente, sino que la estructura de la carga de trabajo cambió. La verdadera enseñanza del curso es esta: al evaluar el hardware de inferencia, no basta con mirar un solo indicador de FLOPs; también hay que revisar si el ancho de banda, la jerarquía de la caché y el patrón de acceso corresponden a las características de uso intensivo de memoria propias de Decode.

\begin{knowledgebox}{Lo que más se pasa por alto al desplegar}
Muchos equipos solo miran el rendimiento promedio del benchmark y pasan por alto la distribución de longitudes de los prompts y el tipo de solicitudes de los usuarios. Si el negocio se basa principalmente en prompts largos con respuestas cortas, Prefill puede ser el cuello de botella principal; si se basa principalmente en prompts cortos con respuestas largas, el cuello de botella de ancho de banda de Decode se vuelve más pronunciado.
\end{knowledgebox}

\subsection{Resumen de la sección}
La diferencia entre Prefill y Decode no es solo conocimiento algorítmico, sino conocimiento de diseño de arquitectura de servicio. Solo al considerar en conjunto las métricas, el hardware y el tipo de carga de trabajo, el ajuste de la inferencia no se desviará del rumbo.

\newpage
